\documentclass{article}

\usepackage{amsmath,amssymb}
\usepackage{xurl}
\usepackage[hidelinks]{hyperref}
\setlength{\emergencystretch}{2em}

% Shared commands for notes in docs/paper-gaps/.

\DeclareUnicodeCharacter{2080}{\ifmmode{}_{0}\else\textsubscript{0}\fi}
\DeclareUnicodeCharacter{2081}{\ifmmode{}_{1}\else\textsubscript{1}\fi}
\DeclareUnicodeCharacter{2082}{\ifmmode{}_{2}\else\textsubscript{2}\fi}
\DeclareUnicodeCharacter{2099}{\ifmmode{}_{n}\else\textsubscript{n}\fi}

\newcommand{\Lean}{\textsc{Lean}}
\ExplSyntaxOn
\NewDocumentCommand{\leanid}{m}{
  \ifmmode
    \textsf{\detokenize{#1}}
  \else
    \group_begin:
    \urlstyle{sf}
    \tl_set:Nn \l_tmpa_tl {#1}
    \tl_replace_all:Nnn \l_tmpa_tl {\_} {_}
    \exp_args:NV \path \l_tmpa_tl
    \group_end:
  \fi
}
\ExplSyntaxOff
\providecommand{\tr}{\operatorname{tr}}
\newcommand{\tnleanIssueBase}{https://github.com/LionSR/TNLean/issues/}
\newcommand{\tnleanPrBase}{https://github.com/LionSR/TNLean/pull/}
\newcommand{\tnleanDocBase}{https://sirui-lu.com/TNLean/docs/}
\newcommand{\ghissue}[1]{\href{\tnleanIssueBase#1}{GitHub issue \##1}}
\newcommand{\ghpr}[1]{\href{\tnleanPrBase#1}{PR \##1}}
\newcommand{\doclink}[1]{\href{\tnleanDocBase#1.html}{\path{#1}}}

% Dirac notation and the identity operator, as in blueprint/src/macros/common.tex.
% \providecommand keeps note-local definitions authoritative.
\usepackage{dsfont}
\providecommand{\ket}[1]{|#1\rangle}
\providecommand{\bra}[1]{\langle #1|}
\providecommand{\braket}[2]{\langle #1 | #2 \rangle}
\providecommand{\ketbra}[2]{|#1\rangle\!\langle #2|}
\providecommand{\Id}{\mathds{1}}
\providecommand{\one}{\mathds{1}}
\providecommand{\C}{\mathbb{C}}
\providecommand{\MN}[1]{M_{#1}(\C)}
\providecommand{\MD}{M_D(\C)}

% External referencing for paper sources.  The build helper writes sanitized
% auxiliary files under build/paper-aux/ and excludes duplicated source labels.
\usepackage{xr}

% Import a generated paper auxiliary file with a caller-supplied label prefix.
% The candidate paths support builds from docs/paper-gaps/, the repository root,
% and build/paper-aux/ itself.
\newcommand{\importpaperaux}[2]{%
  \IfFileExists{../../build/paper-aux/#2.aux}{%
    \externaldocument[#1]{../../build/paper-aux/#2}%
  }{%
    \IfFileExists{build/paper-aux/#2.aux}{%
      \externaldocument[#1]{build/paper-aux/#2}%
    }{%
      \IfFileExists{#2.aux}{%
        \externaldocument[#1]{#2}%
      }{}%
    }%
  }%
}

% General external reference macros
\newcommand{\paperref}[2]{\ref{#1-#2}}
\newcommand{\papereqref}[2]{(\ref{#1-#2})}

% CPSV16 source (1606.00608 / MPDO-22-12-17-2.tex) external document and shortcuts
\importpaperaux{cpsv16-}{1606.00608/MPDO-22-12-17-2-tnlean}
\newcommand{\cpsvref}[1]{\paperref{cpsv16}{#1}}
\newcommand{\cpsveqref}[1]{\papereqref{cpsv16}{#1}}

% Verdict marker: \gapnote{<kind>}{<status>}. Kind is the policy
% classification (clarification, local-correction, scope-restriction,
% unfaithful, false-source, open-gap); status is open, wip (elimination
% underway), resolved, or historical. Machine-read by the site index;
% prints a verdict line.
\newcommand{\gapnote}[2]{\par\noindent\textbf{Verdict:} #1 (#2).\par}


\title{The Domain Boundary of the Topological Gibbs Decomposition}
\author{}
\date{2026-07-23}

\begin{document}
\maketitle
\gapnote{local-correction}{resolved}

\section{Source range}

Theorem~4.15 of Ref.~\cite[lines~999--1016]{Cirac2016MPDO_arXiv} states that
the length-$N$ matrix product density operator has a decomposition
\begin{align}
    \rho^{(N)}(M)
        &= \sum_{i=1}^{d}\lambda_iP_i^{(N)}e^{-H_N},
    \label{eq:cpsv16_topological_gibbs_length_one_source_decomposition}\\
    [P_i^{(N)},e^{-H_N}]
        &= 0.
    \label{eq:cpsv16_topological_gibbs_length_one_source_commutation}
\end{align}
Here
\begin{align}
    H_N
        &= \sum_i h_{i,i+1}
    \label{eq:cpsv16_topological_gibbs_length_one_source_hamiltonian}
\end{align}
is a translationally invariant commuting nearest-neighbor Hamiltonian.

Definition~4.8 of Ref.~\cite[lines~829--847]{Cirac2016MPDO_arXiv} constructs
each local term by placing a two-spin matrix $h$ on the first two sites and
translating it around the periodic chain. This construction is defined for
$N\geq2$. The periodic relation $N+1=1$ identifies site labels; it does not
identify the two tensor factors on which $h$ acts.

\section{The retained-coordinate construction}

Let $\mu$ be the strictly positive diagonal one-site multiplicity matrix in
the retained coordinates. Define the two-site interaction by
\begin{align}
    h
        &= \operatorname{diag}(-\log\mu)\otimes I.
    \label{eq:cpsv16_topological_gibbs_length_one_retained_interaction}
\end{align}
For every $N\geq2$, the periodic translates of $h$ are well-defined,
diagonal, Hermitian, and pairwise commuting. Every site appears exactly once
as the first site of a translated term. Consequently,
\begin{align}
    \exp\left(-\sum_i h_{i,i+1}\right)
        &= \mu^{\otimes N}.
    \label{eq:cpsv16_topological_gibbs_length_one_retained_gibbs_factor}
\end{align}
Combining
Eq.~\ref{eq:cpsv16_topological_gibbs_length_one_retained_gibbs_factor} with
the terminal spectral refinement proves the retained-coordinate Gibbs
decomposition for every $N\geq2$.

Applying the adjoint retained-row map reconstructs the physical density
operator. This reconstruction alone does not provide a physical-space
Hamiltonian or physical-space projectors. The independent physical-coordinate
gap is recorded in
\path{docs/paper-gaps/cpsv16_topological_gibbs_physical_complement.tex}.

The terminal spectral sectors are indexed by pairs $(\gamma,k)$. Positivity of
every multiplicity embeds one such sector into each retained copy, and the
coisometry rank bound gives
\begin{align}
    \#\{(\gamma,k)\}
        &\leq d.
    \label{eq:cpsv16_topological_gibbs_length_one_sector_bound}
\end{align}
Choose an injection of the terminal sectors into
$\{1,\ldots,d\}$, and extend the terminal eigenvalues and projectors by zero.
The added zero matrices are self-adjoint idempotents and commute with the
Gibbs factor. This construction produces exactly $d$ summands without claiming
that there are $d$ nonzero terminal sectors.

\section{The printed domain boundary}

For $N=1$, Definition~4.8 of
Ref.~\cite[lines~829--847]{Cirac2016MPDO_arXiv} defines neither an operator on
the one-site Hilbert space nor a convention for a coincident periodic bond
$h_{1,1}$. In particular, it supplies no map from two-site operators to
one-site operators and no multiplicity for the self-edge.

Several possible extensions would add new data: one could replace $h$ by
$\operatorname{diag}(-\log\mu)$, trace out one tensor factor, restrict to a
diagonal subspace, or otherwise collapse the two tensor factors. None of these
conventions appears in the cited definition.

The source construction therefore determines the theorem only on the domain
\begin{align}
    N
        &\geq 2.
    \label{eq:cpsv16_topological_gibbs_length_one_defined_domain}
\end{align}
Theorem~4.15 of Ref.~\cite[lines~999--1016]{Cirac2016MPDO_arXiv} does not print
this lower bound. Its length-one instance is undefined unless the authors or
an erratum specify both $H_1$ and the coincident-edge convention.

The formal retained- and physical-coordinate theorems therefore make no
length-one assertion. The physical-coordinate complement is complete for
$N\geq2$; see
\path{docs/paper-gaps/cpsv16_topological_gibbs_physical_complement.tex}. That
completion is independent of the absent length-one convention and does not
supply it.

\bibliographystyle{alpha}
\bibliography{references}

\end{document}
